\renewcommand{\thetable}{3.A.14}
\begin{tablalafrox}{Siglas y códigos utilizados. Fuente: terminología de las Bases y de los Subdocumentos 1 y 2.}%
  {tab:3-A-14}%
  {F{0.132} Y}%
  {\cab{Sigla o código} & \cab{Significado}}
  BA / BTT & Bases Administrativas / Bases Técnicas Transversales. \\
  SD1 / SD2 & Subdocumento 1 (Presentación de la empresa) / Subdocumento 2 (Comprensión del problema). \\
  F-01 a F-14 & Familias funcionales del Subdocumento 2, Anexo 2.1. \\
  RF / RNF & Requerimiento funcional / no funcional. El sufijo -BTT distingue un requerimiento de las Bases que comparte número con uno del caso. \\
  RNF-DISP, -REND, -RESIL, -SEG, -DR & Condiciones no funcionales rectoras del Subdocumento 2, Anexo 2.1. \\
  RT & Requisito codificado de las Bases Técnicas Transversales. \\
  R18-01 a R18-16 & Identificadores de los dieciséis resultados de aceptación de Distribuidora Puelche, detallados en el Anexo 3.J. \\
  RNG & Regla de negocio (Anexo 3.G). \\
  S / E / R / D / V & Supuesto / exclusión / restricción / decisión / consulta. \\
  M1 a M12 & Módulos de negocio de la solución (Subdocumento 3, sección 3.4.2). \\
  E1 / E2 & Etapa 1 / Etapa 2 del calendario contractual (\textit{Bases Administrativas}, 2026, art. 17.1, p. 12). \\
  CD & Centro de distribución. \\
  SKU & Código interno que identifica un producto o referencia de inventario (Stock Keeping Unit). \\
  OC & Orden de compra. \\
  TI & Tecnologías de la información. \\
  IA & Inteligencia artificial. \\
  EDT & Estructura de descomposición del trabajo. \\
  ADR & Registro de decisiones de arquitectura (Architecture Decision Record): documenta la decisión, las alternativas y su fundamento. \\
  DEV / QA / PREPROD / PROD / DR & Ambientes de desarrollo / aseguramiento de calidad y pruebas / preproducción / producción / recuperación ante desastres. \\
  ERP & Sistema de gestión empresarial conservado como registro contable y emisor tributario. \\
  WMS & Sistema de gestión de almacenes. \\
  FEFO & Primero en vencer, primero en salir. \\
  OTIF & Entregas completas y a tiempo, según S-01 y RNG-09. \\
  POD & Prueba de entrega. \\
  GDE & Guía de despacho electrónica. \\
  EDI / ASN & Intercambio electrónico estructurado / aviso de despacho. \\
  SSCC & Código seriado de unidad logística GS1. \\
  GS1 & Organización y sistema de estándares para identificar productos y unidades logísticas e intercambiar sus datos. \\
  GTIN & Número global de artículo comercial (Global Trade Item Number), utilizado para identificar productos bajo estándares GS1. \\
  AI & Identificador de aplicación GS1 (Application Identifier): código que indica el significado y formato del dato que lo sigue, como lote o vencimiento. Se distingue de IA (inteligencia artificial). \\
  RTO / RPO & Tiempo máximo de recuperación / pérdida máxima de datos admisible. \\
  NOC & Centro de operaciones de red. \\
  CI/CD & Integración y despliegue continuos. \\
  MFA & Autenticación multifactor: verificación de identidad mediante dos o más factores de distinta naturaleza. \\
  SSO & Inicio de sesión único (Single Sign-On): permite acceder a varios servicios con una misma autenticación. \\
  RBAC & Control de acceso basado en roles (Role-Based Access Control): asigna permisos según el rol del usuario. \\
  TLS & Protocolo de seguridad de la capa de transporte (Transport Layer Security), utilizado para proteger las comunicaciones en tránsito. \\
  SIEM & Gestión de información y eventos de seguridad (Security Information and Event Management): centraliza y correlaciona registros para detectar incidentes. \\
  EDR & Detección y respuesta en equipos finales (Endpoint Detection and Response): supervisa dispositivos y permite responder a amenazas. \\
  SAST & Pruebas estáticas de seguridad de aplicaciones (Static Application Security Testing): analizan el código sin ejecutar la aplicación. \\
  SCA & Análisis de composición de software (Software Composition Analysis): identifica dependencias, vulnerabilidades conocidas y licencias. \\
  DAST & Pruebas dinámicas de seguridad de aplicaciones (Dynamic Application Security Testing): evalúan la aplicación en ejecución. \\
  SBOM & Lista de materiales de software (Software Bill of Materials): inventario de componentes y dependencias del software. \\
\end{tablalafrox}
